\section{Checks for the unequal-coupling estimates}\label{cl:app:anisotropic-checks}

The checker uses only the Python standard library. From the project
repository linked in Data and code availability run
\begin{verbatim}
python scripts/seam_disagreement_constants.py
\end{verbatim}
The computation checks the rational inequalities and sample periodic graphs
below. The polymer expansion and the passage to continuous time are proved in
Theorems~\ref{cl:thm:unequal-gap} and \ref{cl:thm:electric-packet}.

For the $\SU(2)$ parameters in Theorem~\ref{cl:thm:unequal-gap}, put
$u=[2b/(3+2b)]^h$ and $\delta=10^{-6}$. With
$h=\lceil18(1+2b/3)\rceil$ one has $u\le e^{-18}$, also at $b=0$.
The fourth partial sum for $e^3$ is $131/8>16$. Thus
$u<16^{-6}<1/8$, and
\[
 \frac{1+u}{(1-u)^3}-1
 =\frac{4u-3u^2+u^3}{(1-u)^3}
 \le\frac{4u}{(7/8)^3}<6u<\delta.
\]
The inequalities
\[
 \frac{(1+\delta)^4}{1-\delta}<4,\qquad
 4\delta < \frac{1}{2\cdot 3^2\cdot80^2}
 <\frac1{2e^2\,80^2}
\]
certify the remaining activity comparison for $h\lambda\le\delta$.
Here $e^x-1\le x/(1-x)$ for $0\le x<1$, and $e<3$ follows directly
from its factorial series. At $A=a_0$, the ratio in the geometric
sum for polymers containing a given mark is $80^2 A e^2=1/2$.
Its sum is $1/80^2$, which leaves room for the factor $80$ in
counting the closed neighbours of each mark.

The incidence check constructs the columns and plaquettes of the spatial
boxes $(4,4,4)$ and $(4,5,6)$ at $8$ and $12$ time slabs, respectively.
It also constructs both doubled time graphs. Each endpoint belongs to ten
marks, and each mark has two or eight endpoints, according to its type.
The script counts the closed neighbours from these incidences and compares
local patterns rooted across the periodic seam with their two lifts.
The equality of the rooted contributions for general periods follows
from the lifting argument in the theorem. The factor in the raw tail
bound is $2(6Vm)+6V(2m)=24Vm$, since the two periods share the cutoff
$m/4$. Projection then gives the coefficient $48Vm$ for the physical
defect.\footnote{The purity defect need not decrease when one deletes
excited levels. For normalized weights $1$ and thirty copies of $1/10$,
the exponential of the defect is $160/13$. Adding an excited weight $9/10$
makes it $2401/211<160/13$. This finite diagonal example explains the use
of the trace bound in the projection argument.}

For Theorem~\ref{cl:thm:electric-packet}, the upper bound
$12B\kappa<\delta$ implies
$\epsilon^2\le 2B\kappa/3<\delta/18<(1/50)^2$. At $\epsilon=1/50$,
the lower variance bound is $2083/10000>1/5$. The four links of the
plaquette give relative energy at most $4/(2B)$, and division by
$1/5$ gives the Rayleigh bound $10/B$. The relation $th\to12B$
turns $1/(4h)$ into the lower gap $1/(48B)$ in continuous time.
These checks concern the unequal-coupling regime and continuous time on
a fixed spatial lattice.
